% q0068.tex — Foodbook
\question{
  id        = {0068},
  title     = {Foodbook},
  source    = {OBECON},
  year      = {2026},
  round     = {Seletiva IEO -- Phase C},
  language  = {English},
  topic     = {Micro},
  subtopic  = {Information Cascades / Bayesian Updating},
  type      = {dissertative},
  statement = {%
    A restaurant is either Good (probability $\tfrac{1}{2}$) or Bad
    (probability $\tfrac{1}{2}$). Each customer receives an independent
    private signal: ``Good'' with probability $\tfrac{2}{3}$ if the
    restaurant is Good, and ``Bad'' with probability $\tfrac{2}{3}$ if it
    is Bad. Entering gives payoff $+1$ if the restaurant is Good and $-1$
    if Bad; not entering gives payoff $0$.

    \begin{enumerate}[(a)]
      \item (20 points) Alice is first. She receives a ``Bad'' signal.
            Calculate her posterior probability that the restaurant is Good
            and determine her optimal action.

      \item (20 points) Bob is second. He observes that Alice did not enter.
            What does Bob infer about Alice's signal? What is his posterior
            that the restaurant is Good, and what is his optimal action?

      \item (30 points) Charles is third. He observes that both Alice and
            Bob did not enter, and receives a ``Good'' signal. What does
            Charles infer from the history? What is his optimal action?
            Explain how an \emph{information cascade} forms.

      \item (30 points) From the fourth customer onward, does the cascade
            persist even if customers continue receiving ``Good'' signals?
            Discuss the efficiency implications and provide two real-world
            examples of information cascades.
    \end{enumerate}
  },
  choices   = {},
  answer    = {},
  solution  = {%
    \textbf{(a)} Posterior (Good $|$ Bad signal):
    \[
      P(G|\text{Bad}) = \frac{(1/3)(1/2)}{(1/3)(1/2)+(2/3)(1/2)} = \frac{1}{3}.
    \]
    Expected payoff from entering: $(1/3)(+1)+(2/3)(-1) = -1/3 < 0$.
    Alice should \textbf{not enter}.

    \textbf{(b)} Bob knows Alice's not-entering reveals she received a
    ``Bad'' signal. His posterior is also $P(G|\text{Bad signal}) = 1/3$.
    Expected payoff $= -1/3 < 0$. Bob should \textbf{not enter}.

    \textbf{(c)} Charles infers: Alice got Bad ($\Rightarrow$ prior drops to
    $1/3$). If Bob got Good, $P(G) = 1/2$ and he'd be indifferent (enters by
    assumption). Since Bob also didn't enter, Bob got Bad. Prior for Charles:
    $P(G) = 1/5$. Updating on Charles's Good signal:
    \[
      P(G|\text{Good, history}) = \frac{(2/3)(1/5)}{(2/3)(1/5)+(1/3)(4/5)}
      = \frac{1}{3}.
    \]
    Expected payoff $= -1/3 < 0$. Charles should \textbf{not enter} despite
    his favourable private signal. An \textbf{information cascade} has
    formed: the public history overwhelms private signals, and rational
    agents ignore their own information in favour of observed behaviour.

    \textbf{(d)} Yes, the cascade persists: subsequent customers reason the
    same way as Charles, ignoring Good signals. This is \textbf{dynamically
    inefficient}: private information is wasted, and a Good restaurant may
    go patronised only if early customers happened to receive Bad signals.

    \textit{Real-world examples}:
    \begin{itemize}
      \item \textbf{Financial herding}: investors observe others buying or
            selling an asset and follow the crowd, ignoring private
            information --- driving bubbles and crashes.
      \item \textbf{Technology lock-in}: consumers adopt a technology (e.g.\
            VHS over Betamax) because others did, even if an alternative is
            superior, creating welfare-reducing path dependence.
    \end{itemize}
  },
}
